\subsection*{Wykładniki}

\begin{przyklad}
W kategorii $\Set$ mamy: dla wybranych zbiorów $B, C$ istnieje zbiór
\[
C^B \stackrel{df}{=} \{f : B \to C\}.
\]
\end{przyklad}

\begin{uwaga}
Weźmy funkcję $f : A \times B \to C$. Wybierzmy $a \in A$ i zauważmy, że
\[
f_a \stackrel{df}{=} f(a, -) : B \to C
\]
jest elementem zbioru $C^B$. Dla funkcji $f : A \times B \to C$ zdefiniowaliśmy więc funkcję
\[
\hat{f} : A \to C^B, \qquad a \mapsto f_a.
\]
Co więcej, dla każdej funkcji $g : A \to C^B$ istnieje funkcja $f : A \times B \to C$ taka, że $\hat{f} = g$ -- mianowicie $f(x, y) = g(x)(y)$.

Podsumowując, istnieje bijekcja
\[
\Set(A \times B, C) \;\cong\; \Set(A, C^B).
\]
\end{uwaga}

\begin{uwaga}
Bijekcja daje nam przekształcenie
\[
\mathrm{eval} : C^B \times B \to C, \qquad \mathrm{eval}(g, b) = g(b).
\]
Ponadto dla dowolnego $A$ i $f : A \times B \to C$ istnieje jednoznaczne $\tilde{f} : A \to C^B$ takie, że diagram
\[
\begin{tikzcd}
C^B \times B \arrow[r, "\mathrm{eval}"] & C \\
A \times B \arrow[u, "\tilde{f} \times \id_B"] \arrow[ur, "f"'] &
\end{tikzcd}
\]
jest przemienny, tj.\ $\mathrm{eval}\circ(\tilde f\times\id_B) = f$.
\end{uwaga}

\begin{definicja}
Załóżmy, że $\C$ ma binarne produkty. Obiektem potęgowym (\emph{wykładnikiem}) obiektów $B, C \in \C$ nazywamy obiekt $C^B$ wraz ze strzałką
\[
\varepsilon : C^B \times B \to C
\]
taką, że dla każdego obiektu $Z \in \C$ i strzałki $f : Z \times B \to C$ istnieje dokładnie jedna strzałka
\[
\tilde{f} : Z \to C^B
\]
taka, że diagram
\[
\begin{tikzcd}
C^B \times B \arrow[r, "\varepsilon"] & C \\
Z \times B \arrow[u, "\tilde{f} \times \id_B"] \arrow[ur, "f"'] &
\end{tikzcd}
\]
jest przemienny.
\end{definicja}

\begin{uwaga}[Notacja]
\begin{enumerate}
\item $\varepsilon : C^B \times B \to C$ nazywamy \emph{ewaluacją}.
\item $\tilde{f} : Z \to C^B$ nazywamy \emph{transpozycją} (\emph{wykładnikiem}) strzałki $f : Z \times B \to C$.
\end{enumerate}
\end{uwaga}

\begin{uwaga}
Mając daną strzałkę $g : Z \to C^B$, definiujemy
\[
\overline{g} \stackrel{df}{=} \varepsilon \circ (g \times \id_B) : Z \times B \to C,
\]
tj.
\[
\begin{tikzcd}
Z \times B \arrow[r, "g \times \id_B"] & C^B \times B \arrow[r, "\varepsilon"] & C
\end{tikzcd}
\]
Wtedy $\overline{\tilde{f}} = f$ oraz $\widetilde{\overline{g}}=g$. Innymi słowy, jeśli kategoria $\C$ ma obiekty potęgowe, to mamy bijekcję
\[
\C(Z \times B, C) \;\cong\; \C(Z, C^B), \qquad f \mapsto \tilde{f}, \quad g \mapsto \overline{g}.
\]
\end{uwaga}

\subsection*{Kategorie kartezjańsko domknięte}

\begin{definicja}
Kategorię $\C$ nazywamy \emph{kartezjańsko domkniętą}, jeśli ma:
\begin{enumerate}
\item wszystkie skończone produkty,
\item wszystkie obiekty potęgowe (dla dowolnej pary $B, C$).
\end{enumerate}
\end{definicja}

\begin{przyklad}
$\Set$ oraz $\Pos$ są kategoriami kartezjańsko domkniętymi.
\end{przyklad}

\begin{stwierdzenie}
$\Pos$ jest kartezjańsko domknięta.
\end{stwierdzenie}

\begin{proof}[Szkic]
\textbf{(1)} $\Pos$ ma wszystkie skończone produkty: jednoelementowy poset jest terminalny w $\Pos$, a dla posetów $P = (P, \leq_P)$, $Q = (Q, \leq_Q)$ definiujemy
\[
P \times Q = (P \times Q, \leq),
\]
gdzie
\[
(p_1, q_1) \leq (p_2, q_2) \;\stackrel{df}{\iff}\; p_1 \leq_P p_2 \;\wedge\; q_1 \leq_Q q_2.
\]
Razem z rzutowaniami $\pi_1, \pi_2$ daje to produkt.

\textbf{(2)} Dla posetów $P$ i $Q$ definiujemy
\[
Q^P \stackrel{df}{=} \big(\{f : P \to Q \mid f \text{ monotoniczna}\},\ \leq\big),
\]
gdzie
\[
f_1 \leq f_2 \;\stackrel{df}{\iff}\; f_1(p) \leq_Q f_2(p) \quad \text{dla każdego } p \in P,
\]
oraz ewaluację
\[
\varepsilon : Q^P \times P \to Q, \qquad \varepsilon(f, p) \stackrel{df}{=} f(p),
\]
która jest dobrze zdefiniowaną strzałką w $\Pos$. Transpozycja $\tilde{f}\colon X\to Q^P$ jest zdefiniowana analogicznie jak w $\Set$.
\end{proof}
